\section{Corredor excepcional, teoría conforme, cuerdas y clasificación de masas}
\Needspace{7\baselineskip}
\subsection{Preimagen común y coordinación de las \(2\) lecturas del corredor excepcional}
La relación entre masas y simetrías excepcionales no comienza en una acción del Monstruo sobre una tabla de partículas. Comienza en el mismo estado enriquecido del TPK. Sea

\[
x\in\mathcal X_{\rm TPK}^{\rm enr}.
\]
La lectura arquimediana contrae cilindros compatibles y conserva el valor real con su genealogía:

\[
\mathcal P_{\rm arq}(x)\in\mathbb R.
\]
La lectura excepcional conserva incidencias de frontera y produce una clase en el límite proyectivo:

\[
\mathcal P_{\rm exc}(x)=Q_\infty(x).
\]
La naturalidad exige, para todo refinamiento \(n\ge m\),

\[
r_{n,m}\circ Q_n=Q_m\circ p_{n,m}.
\]
La aplicación conjunta

\[
\mathcal P_{\rm joint}:x\longmapsto
\bigl(\mathcal P_{\rm arq}(x),\mathcal P_{\rm exc}(x)\bigr)
\]
es la forma positiva de la doble proyección. No niega la conexión Monstruo--masa: la tipa. La rama de masa y la rama excepcional son imágenes de una misma preimagen con orientación, gradación, acarreo, memoria y holonomía comunes. Tampoco afirma que el Monstruo actúe literalmente sobre las cifras decimales o que un coeficiente de Moonshine sea por sí mismo una masa.

\Needspace{7\baselineskip}
\subsection{Construcción del corredor excepcional}
La incidencia ternaria prolonga el estado por la cadena

\[
C_3
\longrightarrow A_2^{12}
\longrightarrow N(A_2^{12})
\longrightarrow \Lambda_{24}
\longrightarrow V^\natural
\longrightarrow \mathbb M,
\]
donde \(\Lambda_{24}\) es la red de Leech, \(V^\natural\) el álgebra de vértices del Moonshine y \(\mathbb M\) el grupo Monstruo. El paso no consiste en yuxtaponer nombres clásicos: cada flecha debe preservar la incidencia, la gradación y la compatibilidad de refinamiento heredadas del estado.

La traza ternaria se escribe

\[
t_3(q)=q^{-1}
\prod_{n\ge1}(1+q^n+q^{2n})^{-12}.
\]
El primer defecto no trivial satisface

\[
D_{\rm APP}=[q]t_3=v_\alpha^2=54,
\]
y el primer coeficiente del carácter del Monstruo conserva la descomposición

\[
196884=54+10\cdot3^9.
\]
Estas identidades transportan gradación y multiplicidad. No asignan una masa en MeV. Su función en la teoría de masas es clasificar el espacio sobre el que actúa el operador antes de escalarizarlo.

\Needspace{7\baselineskip}
\subsection{Emparejamiento fibra a fibra con la ley de masas}
La rama material del mismo estado produce

\[
x\longmapsto
\bigl(F_X,\sigma_X,\eta_X,\widehat M_X\bigr),
\]
donde \(F_X\) es punto, órbita, fibra o multisección; \(\sigma_X\in\mathbb Z^5\) es la firma; \(\eta_X\in\mathcal E_{90,120}\) es el refinamiento; y \(\widehat M_X\) el operador de masa. La lectura conjunta es entonces

\[
\mathfrak J_{\rm mass,exc}(x)
=\bigl(F_X,\sigma_X,\eta_X,\widehat M_X,Q_\infty(x)\bigr).
\]
Dos estados con el mismo valor escalar de masa pueden permanecer distintos por su componente excepcional; dos estados con el mismo carácter excepcional pueden poseer operadores de masa diferentes. La clasificación correcta es el emparejamiento de fibras, no la igualdad aislada de números.

Si una representación excepcional \(G\) actúa en una fibra física \(\mathcal H_X\), la compatibilidad con la masa exige

\[
U_g\widehat M_XU_g^{-1}=\widehat M_X,
\qquad g\in G.
\]
La descomposición

\[
\mathcal H_X=\bigoplus_{\lambda}
\mathcal H_{X,\lambda}
\]
separa multipletes excepcionales. Sobre un sumando irreducible, la conmutación fuerza que la compresión de \(\widehat M_X\) sea escalar. Ésta es una vía legítima de escalarización por simetría. El valor del escalar sigue procediendo del carácter, las torres, la sección dimensional y el acoplamiento; la dimensión de la representación sólo determina la degeneración.

\Needspace{7\baselineskip}
\subsection{Fock y álgebra de vértices}
La prolongación de una ruta a un espacio de muchas partículas posee dos completaciones principales. La exterior produce CAR, Pauli y estadística fermiónica; la simétrica produce CCR y estadística bosónica. Una álgebra de vértices añade producto operatorio, vacío, traslación y gradación:

\[
V=\bigoplus_{n\ge0}V_n,
\qquad
L_0|_{V_n}=nI.
\]
El carácter

\[
\operatorname{ch}_V(q)
=\operatorname{Tr}_V q^{L_0-c/24}
\]
cuenta estados por grado. Una multiplicidad en \(V_n\) no es una masa. Para convertir la gradación en energía se necesita un Hamiltoniano y una sección de longitud o tiempo. En una realización conforme sobre un cilindro de circunferencia \(L\),

\[
H_{\rm CFT}
=\frac{2\pi\hbar_{\rm int}v}{L}
\left(L_0+\bar L_0-\frac{c}{12}\right).
\]
La masa o brecha física aparece sólo después de fijar \(v\), \(L\), las condiciones de borde y el mapa desde la ruta TPK a los módulos conformes. La función del corredor excepcional consiste en organizar las degeneraciones y las reglas de composición de esos módulos.

\Needspace{7\baselineskip}
\subsection{Espectro de cuerda como realización dimensional posterior}
Sea \(\alpha'\in\mathcal L_L^2\) la longitud cuadrática que representa la tensión inversa en la carta naturalizada de cuerda declarada. En una compactificación circular de radio \(R\), los momentos izquierdo y derecho son

\[
p_L=\frac{m}{R}+\frac{wR}{\alpha'},
\qquad
p_R=\frac{m}{R}-\frac{wR}{\alpha'},
\]
con \(m,w\in\mathbb Z\). En unidades naturales de la realización, el espectro cerrado satisface

\[
M^2=p_L^2+\frac{4}{\alpha'}(N_L-a_L)
=p_R^2+\frac{4}{\alpha'}(N_R-a_R),
\]
junto con la condición de nivel

\[
N_L-N_R=a_L-a_R-mw.
\]
La involución

\[
(m,w,R)\longleftrightarrow
\left(w,m,\frac{\alpha'}{R}\right)
\]
conserva el espectro: es la realización de dualidad T. HMT aporta antes la orientación bidireccional, la memoria y el emparejamiento de hojas; la teoría de cuerda añade tensión, radio, osciladores, nivel y condiciones de contorno. No se obtiene una masa de cuerda reemplazando \(N_L,N_R\) por alturas de la torre \(\langle90,120\rangle\). Toda identificación entre ambas gradaciones requiere un entrelazador que conserve suma, orientación y carácter.

\Needspace{7\baselineskip}
\subsection{Pantallas reticulares de teoría M: gradación, compactificación y transporte de la memoria residual}
Las pantallas internas

\[
12\longrightarrow11\longrightarrow10,
\qquad
26\longrightarrow24,
\]
conservan el descenso dodecafásico, la eliminación de una dirección redundante y la pantalla transversal. En una realización de teoría M, los campos

\[
g_{MN},\qquad C_{MNP},\qquad\Psi_M
\]
deben obtenerse mediante un entrelazador que preserve calibre, supersimetría, restricciones y producto interno. Las pantallas numéricas no sustituyen ese mapa.

La clasificación de masas se transporta a esta rama mediante

\[
\text{ruta}
\longrightarrow\text{registro genealógico}
\longrightarrow\text{firma}
\longrightarrow\text{carácter}
\longrightarrow\text{operador}
\longrightarrow\text{espectro compactificado}.
\]
Los modos de Kaluza--Klein, devanado, membrana o excitación oscilatoria son estados distintos aunque compartan una masa. El gravitón de una realización de cuerda cerrada, si aparece, es un modo colectivo sin masa del sector geométrico; no se convierte por ello en semilla elemental del TPK ni invalida la formulación de la gravedad como holonomía y torsión.

\Needspace{7\baselineskip}
\subsection{Realizaciones ordinarias e hipotéticas inducidas por el corredor excepcional}
La rama excepcional permite distinguir cuatro fuentes de multiplicidad:

\begin{enumerate}
\item multiplicidad celular o de ruta en el atlas;
\item degeneración de una representación excepcional;
\item ocupación de Fock;
\item nivel conforme, oscilatorio o compactificado.
\end{enumerate}
Confundirlas produce falsos compañeros, falsas generaciones o masas duplicadas. Un estado gemelo, supersimétrico, oscuro o exótico sólo constituye una nueva especie cuando su quíntuplo

\[
(F_X,Q_\infty(x),\rho_X,\widehat M_X,P_X)
\]
difiere físicamente y el Hamiltoniano posee un polo o una sección persistente adicional. Una coincidencia de dimensión, carácter o masa no basta.

En sentido inverso, una degeneración protegida por el corredor excepcional es una predicción estructural: perturbaciones que preserven la simetría no pueden escindir el multiplete. Una escisión observada exige identificar la ruptura de simetría, el acoplamiento o la carta que la produce.

\Needspace{7\baselineskip}
\subsection{Condiciones de positividad, cierre y estabilidad de las realizaciones topológicas}
La conexión entre el corredor excepcional y las masas queda falsada en una realización concreta si ocurre cualquiera de los hechos siguientes:

\begin{itemize}
\item falla la naturalidad \(r_{n,m}Q_n=Q_mp_{n,m}\);
\item la etiqueta excepcional cambia al permutar las masas externas;
\item una multiplicidad de Moonshine se usa como valor de masa sin Hamiltoniano;
\item el operador de masa no conmuta con la simetría que se afirma preservada;
\item se identifica una altura de torre con un nivel de cuerda sin entrelazador;
\item falla la condición de nivel o la invariancia por dualidad T;
\item una pantalla dimensional se declara teoría M sin construir sus campos;
\item se usa la existencia de un modo de espín dos para redefinir la gravedad
interna anterior.

\end{itemize}
Estos controles mantienen positiva la conexión: existe un estado común y una clasificación conjunta; lo que se prohíbe es sustituir las flechas por una analogía numérica.

\Needspace{7\baselineskip}
\subsection{Construcción matemática de las clases materiales a partir de rutas, firmas y representaciones}
La doble realización, la naturalidad proyectiva, la cadena \(A_2^{12}\)--Leech--\(V^\natural\)--Monstruo, la traza \(t_3\), las identidades \(D_{\rm APP}=54\) y \(196884=54+10\cdot3^9\), y las pantallas dimensionales pertenecen a la arquitectura matemática APP--TRIT--TPK desarrollada previamente. La formulación de \(\mathfrak J_{\rm mass,exc}\), la separación de las cuatro multiplicidades y los controles anteriores hacen explícita su función dentro de una teoría autónoma de masas. Las fórmulas conforme y de cuerda son realizaciones dimensionales posteriores con sus parámetros declarados; no se emplean para generar retroactivamente APP, TRIT, TPK, el carácter ni las torres.
